\chapter{Concurrency}
\label{ch:concurrency}

\haira{} provides Go-style concurrency primitives: lightweight spawned tasks and channels for communication. The philosophy is ``share memory by communicating'' rather than ``communicate by sharing memory.''

\section{Concurrency Model}

\begin{itemize}
    \item \textbf{Spawned tasks} -- Lightweight concurrent execution units
    \item \textbf{Channels} -- Typed conduits for communication
    \item \textbf{Select} -- Multiplexing channel operations
    \item \textbf{No shared mutable state} -- Communication via channels
\end{itemize}

\section{Spawn}

The \keyword{spawn} keyword starts a new concurrent task:

\begin{lstlisting}
import "io"
import "time"

fn main() {
    spawn {
        time.sleep(1000)
        io.println("Hello from spawned task!")
    }

    io.println("Main continues immediately")
    time.sleep(2000)  // Wait for spawned task
}
\end{lstlisting}

\subsection{Spawning Functions}

\begin{lstlisting}
fn worker(id: int) {
    io.println("Worker " + to_string(id) + " started")
    time.sleep(1000)
    io.println("Worker " + to_string(id) + " done")
}

fn main() {
    for i in 0..5 {
        spawn worker(i)
    }
    time.sleep(3000)  // Wait for all workers
}
\end{lstlisting}

\subsection{Spawn Returns Handle}

\begin{lstlisting}
fn compute() -> int {
    time.sleep(1000)
    return 42
}

fn main() {
    handle = spawn compute()

    // Do other work...

    result = await handle  // Wait for result
    io.println("Result: " + to_string(result))
}
\end{lstlisting}

\section{Channels}

Channels are typed conduits for passing values between tasks:

\begin{lstlisting}
// Create a channel
ch = chan<int>()           // Unbuffered channel
ch = chan<int>(10)         // Buffered channel (capacity 10)
ch = chan<string>()        // String channel
\end{lstlisting}

\subsection{Sending and Receiving}

\begin{lstlisting}
ch = chan<int>()

// Send (blocks until received for unbuffered)
ch <- 42

// Receive (blocks until value available)
value = <-ch
\end{lstlisting}

\subsection{Basic Channel Example}

\begin{lstlisting}
import "io"

fn main() {
    ch = chan<string>()

    spawn {
        ch <- "Hello from spawned task!"
    }

    message = <-ch
    io.println(message)
}
\end{lstlisting}

\subsection{Buffered Channels}

\begin{lstlisting}
ch = chan<int>(3)  // Buffer size 3

// These don't block (buffer has space)
ch <- 1
ch <- 2
ch <- 3

// This would block (buffer full)
// ch <- 4

// Receive
io.println(<-ch)  // 1
io.println(<-ch)  // 2
io.println(<-ch)  // 3
\end{lstlisting}

\subsection{Closing Channels}

\begin{lstlisting}
ch = chan<int>(10)

spawn {
    for i in 0..10 {
        ch <- i
    }
    close(ch)  // Signal no more values
}

// Iterate until channel closed
for value in ch {
    io.println(value)
}

io.println("Channel closed, loop exited")
\end{lstlisting}

\subsection{Checking Channel State}

\begin{lstlisting}
// Receive with closed check
value, ok = <-ch
if not ok {
    io.println("Channel closed")
}

// Non-blocking receive
value, ok = ch.try_receive()
if ok {
    io.println("Got: " + to_string(value))
} else {
    io.println("No value available")
}

// Non-blocking send
ok = ch.try_send(42)
if not ok {
    io.println("Channel full or closed")
}
\end{lstlisting}

\section{Select}

The \keyword{select} statement waits on multiple channel operations:

\begin{lstlisting}
import "io"
import "time"

fn main() {
    ch1 = chan<string>()
    ch2 = chan<string>()

    spawn {
        time.sleep(1000)
        ch1 <- "from channel 1"
    }

    spawn {
        time.sleep(500)
        ch2 <- "from channel 2"
    }

    // Wait for first available
    select {
        msg = <-ch1 => io.println("Received: " + msg)
        msg = <-ch2 => io.println("Received: " + msg)
    }
}
\end{lstlisting}

\subsection{Select with Default}

Non-blocking select:

\begin{lstlisting}
select {
    value = <-ch => io.println("Got: " + to_string(value))
    default => io.println("No value ready")
}
\end{lstlisting}

\subsection{Select with Timeout}

\begin{lstlisting}
timeout = time.after(5000)  // 5 seconds

select {
    value = <-ch => io.println("Got: " + to_string(value))
    <-timeout => io.println("Timed out!")
}
\end{lstlisting}

\subsection{Select in Loop}

\begin{lstlisting}
fn worker(jobs: chan<Job>, results: chan<Result>, done: chan<bool>) {
    while true {
        select {
            job = <-jobs => {
                result = process(job)
                results <- result
            }
            <-done => {
                io.println("Worker shutting down")
                return
            }
        }
    }
}
\end{lstlisting}

\section{Common Patterns}

\subsection{Worker Pool}

\begin{lstlisting}
fn worker(id: int, jobs: chan<int>, results: chan<int>) {
    for job in jobs {
        io.println("Worker " + to_string(id) + " processing " + to_string(job))
        time.sleep(100)
        results <- job * 2
    }
}

fn main() {
    num_jobs = 10
    num_workers = 3

    jobs = chan<int>(num_jobs)
    results = chan<int>(num_jobs)

    // Start workers
    for i in 0..num_workers {
        spawn worker(i, jobs, results)
    }

    // Send jobs
    for j in 0..num_jobs {
        jobs <- j
    }
    close(jobs)

    // Collect results
    for _ in 0..num_jobs {
        result = <-results
        io.println("Result: " + to_string(result))
    }
}
\end{lstlisting}

\subsection{Fan-Out, Fan-In}

\begin{lstlisting}
fn producer(out: chan<int>) {
    for i in 0..100 {
        out <- i
    }
    close(out)
}

fn worker(in: chan<int>, out: chan<int>) {
    for value in in {
        out <- value * value
    }
}

fn main() {
    input = chan<int>(100)
    output = chan<int>(100)

    // Producer
    spawn producer(input)

    // Fan-out to multiple workers
    for _ in 0..4 {
        spawn worker(input, output)
    }

    // Fan-in (collect results)
    spawn {
        time.sleep(1000)
        close(output)
    }

    for result in output {
        io.println(result)
    }
}
\end{lstlisting}

\subsection{Pipeline}

\begin{lstlisting}
fn generate(out: chan<int>) {
    for i in 2..100 {
        out <- i
    }
    close(out)
}

fn filter(in: chan<int>, out: chan<int>, prime: int) {
    for n in in {
        if n % prime != 0 {
            out <- n
        }
    }
    close(out)
}

fn sieve() -> chan<int> {
    primes = chan<int>()

    spawn {
        ch = chan<int>(100)
        spawn generate(ch)

        while true {
            prime, ok = <-ch
            if not ok {
                break
            }
            primes <- prime

            next = chan<int>(100)
            spawn filter(ch, next, prime)
            ch = next
        }
        close(primes)
    }

    return primes
}

fn main() {
    for prime in sieve() {
        io.println(prime)
        if prime > 50 {
            break
        }
    }
}
\end{lstlisting}

\subsection{Semaphore}

\begin{lstlisting}
struct Semaphore {
    ch: chan<bool>
}

fn new_semaphore(n: int) -> Semaphore {
    sem = Semaphore{ch: chan<bool>(n)}
    for _ in 0..n {
        sem.ch <- true
    }
    return sem
}

Semaphore.acquire() {
    <-self.ch
}

Semaphore.release() {
    self.ch <- true
}

// Usage: limit concurrent operations
fn main() {
    sem = new_semaphore(3)  // Max 3 concurrent

    for i in 0..10 {
        spawn {
            sem.acquire()
            defer sem.release()

            io.println("Task " + to_string(i) + " running")
            time.sleep(1000)
        }
    }

    time.sleep(5000)
}
\end{lstlisting}

\section{Channel Types}

\subsection{Directional Channels}

Functions can restrict channel direction:

\begin{lstlisting}
// Send-only channel
fn producer(out: chan<-int) {
    out <- 42
    // <-out  // Error: cannot receive from send-only channel
}

// Receive-only channel
fn consumer(in: <-chan<int>) {
    value = <-in
    // in <- 42  // Error: cannot send to receive-only channel
}

fn main() {
    ch = chan<int>()
    spawn producer(ch)   // Converts to chan<-int
    spawn consumer(ch)   // Converts to <-chan<int>
}
\end{lstlisting}

\section{Synchronization}

\subsection{WaitGroup}

Wait for multiple tasks to complete:

\begin{lstlisting}
import "sync"

fn main() {
    wg = sync.WaitGroup{}

    for i in 0..5 {
        wg.add(1)
        spawn {
            defer wg.done()
            io.println("Task " + to_string(i))
            time.sleep(1000)
        }
    }

    wg.wait()  // Block until all done
    io.println("All tasks completed")
}
\end{lstlisting}

\subsection{Mutex}

For rare cases requiring shared state:

\begin{lstlisting}
import "sync"

struct Counter {
    mu: sync.Mutex
    value: int
}

Counter.increment() {
    self.mu.lock()
    defer self.mu.unlock()
    self.value += 1
}

Counter.get() -> int {
    self.mu.lock()
    defer self.mu.unlock()
    return self.value
}
\end{lstlisting}

\begin{warning}
Prefer channels over mutexes. Mutexes are provided for interoperability and performance-critical sections, but channel-based designs are generally safer and more idiomatic.
\end{warning}

\section{Concurrency Grammar}

\begin{lstlisting}[language=EBNF]
spawn_expr = "spawn" ( block | function_call ) ;

channel_type = "chan" "<" type ">"
             | "chan<-" type         (* send-only *)
             | "<-chan<" type ">"    (* receive-only *) ;

channel_create = "chan" "<" type ">" "(" [ expression ] ")" ;

send_stmt = expression "<-" expression ;

receive_expr = "<-" expression ;

select_stmt = "select" "{" { select_case } [ default_case ] "}" ;

select_case = pattern "=" receive_expr "=>" ( expression | block )
            | send_stmt "=>" ( expression | block ) ;

default_case = "default" "=>" ( expression | block ) ;
\end{lstlisting}
